\documentclass[11pt]{scrartcl}

\usepackage[top=1.5cm]{geometry}
\usepackage{url}
\usepackage{float}
\usepackage{listings}
\usepackage{xcolor}
\usepackage{graphicx}
\usepackage{booktabs}
\usepackage{makecell}

\setlength{\parindent}{0em}
\setlength{\parskip}{0.5em}

\newcommand{\youranswerhere}{[Your answer goes here \ldots]}
\renewcommand{\thesubsection}{\arabic{subsection}}

\lstdefinestyle{dmrsql}{
  language=SQL,
  basicstyle=\small\ttfamily,
  keywordstyle=\color{magenta!75!black},
  stringstyle=\color{green!50!black},
  showspaces=false,
  showstringspaces=false,
  commentstyle=\color{gray}}

\lstdefinestyle{dmrJava}{
  language=JAVA,
  basicstyle=\small\ttfamily,
  keywordstyle=\color{magenta!75!black},
  stringstyle=\color{green!50!black},
  showspaces=false,
  showstringspaces=false,
  commentstyle=\color{gray}}

\title{
  \textbf{\large Projektaufgabe 3 } \\
  Phase 3 – Implementierung des XPath Accelerators (3.5 P) \\
  {\large Datenmanagement jenseits von Relationen}
}

\author{
  Gruppen Nummer 12 \\
  \large Zanotti Christian, 12528509 \\
  \large Wiesinger Martin, 11935667 \\
  \large Ottino David, 51841010 
}

\begin{document}

\maketitle\thispagestyle{empty}

Dieses Reporting Template dient der Vorbereitung der Abgabe von Phase 3. 

\subsection*{Verkleinern der Fensteranfrage (0.5 P)}

Zeigen Sie, dass die Berechnung der Achsen für die Beispiele aus Phase 1 hier die gleichen Ergebnisse erzeugen:

\begin{table}[h]
	\centering
		\begin{center}
			\begin{tabular}{ l | c c }
				\toprule
				Achse & Ergebnisknoten ID's & Ergebnisgröße\\
				\midrule
				ancestor & 0, 1, 2 & 3 \\ 
				descendants &
\makecell[l]{3, 4, 5, 6, 7, 8, 9, 10, 11, 12, 13, 14, 15, \\
16, 17, 18, 19, 20, 21, 22, 23, 24, 25, \\
26, 27, 28, 29, 30}
& 28 \\  
				following SchmittKAMM23& 18 & 1 \\  
				preceeding SchmittKAMM23 & - & 0 \\ 
				following SchalerHS23& - & 0 \\  
				preceeding SchalerHS23& 3 & 1 \\ 
				\bottomrule
			\end{tabular}
			\end{center}
	\caption{Ergebnisse der XPath-Berechnung auf Edge Modell aus Phase 1}
	\label{tab:ErgebnisseDerXPathBerechnug}
\end{table}

\begin{table}[h]
	\centering
		\begin{center}
			\begin{tabular}{ l | c c }
				\toprule
				Achse & Ergebnisknoten ID's & Ergebnisgröße\\
				\midrule
				ancestor & 0, 1, 2 & 3 \\ 
				descendants &
\makecell[l]{3, 4, 5, 6, 7, 8, 9, 10, 11, 12, 13, 14, 15, \\
16, 17, 18, 19, 20, 21, 22, 23, 24, 25, \\
26, 27, 28, 29, 30}
& 28 \\  
				following SchmittKAMM23& 18 & 1 \\  
				preceeding SchmittKAMM23 & - & 0 \\ 
				following SchalerHS23& - & 0 \\  
				preceeding SchalerHS23& 3 & 1 \\ 
				\bottomrule
			\end{tabular}
			\end{center}
	\caption{Ergebnisse der XPath-Berechnung des XPath-Accelerators aus Phase 3 mit Verkleinerung des Fensters}
	\label{tab:ErgebnisseDerXPathBerechnug}
\end{table}



\subsection*{Schemaerstellung (1 Punkt).}
Geben Sie die Create Table statements aller von Ihnen angelegten Tabellen an:
\begin{lstlisting}[style=dmrsql]
[CREATE TABLE accel_single (
                pre_min INT PRIMARY KEY,
                pre_max INT NOT NULL,
                parent INT,
                node_id INT NOT NULL
            );

CREATE TABLE content (
                node_id INT PRIMARY KEY,
                tag TEXT NOT NULL,
                text TEXT
            );

CREATE TABLE attribute (
                node_id INT NOT NULL,
                name TEXT NOT NULL,
                value TEXT NOT NULL,
                PRIMARY KEY (node_id, name)
            );

CREATE INDEX idx_accel_single_pre_min ON accel_single(pre_min);

CREATE INDEX idx_accel_single_pre_max ON accel_single(pre_max);

CREATE INDEX idx_accel_single_parent ON accel_single(parent);

CREATE INDEX idx_content_tag ON content(tag);

CREATE TABLE accel (
                pre INT PRIMARY KEY,
                post INT NOT NULL,
                parent INT,
                node_id INT NOT NULL,
                height INT NOT NULL
            )]
\end{lstlisting}

\subsection*{Zugriff mit nur einer Achse (1 Punkt).}
Zeigen Sie, dass die Berechnung der Achsen für die Beispiele aus Phase 1 hier die gleichen Ergebnisse erzeugen:

\begin{table}[h]
	\centering
		\begin{center}
			\begin{tabular}{ l | c c }
				\toprule
				Achse & Ergebnisknoten ID's & Ergebnisgröße\\
				\midrule
				Phase 1 &\makecell[l]{3, 4, 5, 6, 7, 8, 9, 10, 11, 12, 13, 14, 15, \\
16, 17, 18, 19, 20, 21, 22, 23, 24, 25, \\
26, 27, 28, 29, 30}
& 28 \\  
Phase 3 &\makecell[l]{3, 4, 5, 6, 7, 8, 9, 10, 11, 12, 13, 14, 15, \\
16, 17, 18, 19, 20, 21, 22, 23, 24, 25, \\
26, 27, 28, 29, 30}
& 28 \\ 
				\bottomrule
			\end{tabular}
			\end{center}
	\caption{Ergebnisse der XPath-Berechnung des XPath-Accelerators aus Phase 3 mit nur einer Achse}
	\label{tab:ErgebnisseDerXPathBerechnug}
\end{table}

Da sich die Knoten IDs unterscheiden, zeigen Sie für einige Knoten die jeweilige Identität.

Die Knoten-Ids sind in unserer Implementierung identisch, da in beiden Phasen der gleiche Tree verwendet wird und somit die gleichen IDs vergeben werden.


\subsection*{Zeitmamagement}

Benötigte Zeit pro Person (nur Phase 3): \textbf{5}


\end{document}
